\subsection{1ST ATTEMPT: Values of functions and of families of proper classes, unified approach}
SKIP THIS, GO TO THE NEXT SUBSECTION
\begin{definition}(Value of relation/function/protected family at x)
\[
F(x):=\union\{b|(x,b)\in F\}
\]
\end{definition}
\begin{definition}(Value of a simple family at x)
\[
F(x) :=\{b|(x,b)\in F\}=\ran(F)
\]
(if domain is defined as usual, then the element of a simple family is intrinsicallly non-empty. if domain set to be whole class of sets, then the element of such "global simple family" can be empty. Both assumtions are weird.)
%F(x) :=\union\{b|(x,b)\in F\}
\end{definition}
We need to improve the previous definitions, because:
1) we want to have family with empty representative, we need to define.
2) it if not possible to preserve domain if element of a family turned out to be empty.
\begin{definition}(Value of a protected family at x)
\[
F(x)_p := \{b|(x,\{b\})\in F\}
F(x)_p := \union\{s|(x,s)\in F\}
\]
The latter definition therefore is universal for any of r/f/F/PF
\end{definition}


\begin{definition}(Domain of relation/function/family/protected family)
\[
\Dom(F):=\{x|\exists s:(x,s)\in F\}=:\dom(F)
\]
\end{definition}

\begin{definition}(Range of relation)(!!!Note that it is in general different from range of a function or family. "R-Ran","ran","small range")
\[
\ran(R):=\{y|\exists x:(x,y)\in R\}
\]
\end{definition}

\begin{definition}(In range of relation/function/family. (It works for both R-Ran and F-Ran.))
\[
(K\mbox{ in classrange of }F) \IFF \exists x: (x \in \Dom(F)) \land (F(x)=K)
\]
\end{definition}
\begin{definition}(In range of protected family. (It works for both R-Ran and F-Ran.))
\[
(K\mbox{ in classrange of protected }F) \IFF \exists x: (x \in \Dom(F)) \land (F(x)_p=K)
\]
\end{definition}


\begin{definition}([Set-]Range of function/family. "F-Ran","Ran", "big range")
%\[
%\Ran(F):=\{y|\exists x: (x \in \Dom(F)) \land (F(x)=y)\}
%\]
\[
\Ran(F):=\{y| y\mbox{ in classrange of } F\}
\]
\end{definition}

\begin{definition}(Function)
A function is a relation which is univocal.
\end{definition}
\begin{definition}(Protected Family)
A protected family is a relation which range of relation consists only of subsets of singletones.
\end{definition}


%A relation is a function when it is univocal.
%A relation is a family when it's relation range consists only of subsets of singletones.

\begin{definition}(r/f/F is defined at x)
\[
(F \ldefined\mbox{ at }x) \IFF (x\in\dom(F))
\]
\end{definition}
%For every axiomatic definition, there are certain prerequisites.
It doesn't make sense to prove anything about $(\varepsilon x:\varphi)$, if it is not proven that $\exists x: \varphi$.\\
It doesn't make sense to prove anything about $(\iota x:\varphi)$, if it is not proven that $\exists! x: \varphi$.\\
There is a "default value" D1 which $(\varepsilon x:\varphi)$ is equal to, if it is not true that $\exists x: \varphi$.\\
There is a "default value" D2 which $(\iota x:\varphi)$ is equal to, if it is not true that $\exists x: \varphi$.\\
They are not specified. (Up to now).\\
If they are not specified, then we cannot define the "(- defined)" predicate.
It is because if we have take a D1 or D2 just as classes in some other application, then they will be "undefined", while they are perfectly okay.(for example, the empty set of the class of all sets)\\
It is not possible to have a "defined" predicate, because it isn't preserved under equality.\\
So naive "defined" requires a recursion on a term.\\
Even if we introduce own "defined" for each constructor of the term, it still will require syntactic analysis.\\
So the solution is to have a context under which definitions make sense.\\
It may be a family of classes, where the first one should be non-empty, as a signal that some conditions are satisfied.\\
It strongly reminds a notion of signature and structure.\\

In the situation, when $\exists x: \varphi$ and $\neg\exists! x: \varphi$.
"value-defined".

Relations are symmetrical, so inverse relation is defined.

Protected families are assymetrical, so it is not always possible to define the inverse of a family.

Families 
\subsection{Non-empty families}
All definitions in this section are local: they will not be used in other sections.
\begin{definition}(family of non-empty classes)
\[
(F is NEF) 
\]
\end{definition}



\subsection{Unified approach to functions and of families}

\begin{definition}(Conglomerate of all classes)
This symbol is not a term of our theory(MK+axGroth), because terms denote classes, and it is bigger than a class.
\end{definition}
Different variant of canonisations:
\begin{definition}(Canonization a family of non-empty classes)\\
$
\canonize0(F) = \{(x,\{q\})|\exists y:((x,y)\in F)\land (q\in y)\}
$\\

\end{definition}
\begin{definition}(Canonization[of a family of potentially empty classes])\\
$
\canonize1(F) = \{(x,\{q\})|\exists y:((x,y)\in F)\land (q\in y)\}\cup\{(x,\emptyset)|x\in\Dom(F)\}
$\\
$
\canonize2(F) = \{(x,\{q\})|\exists y:((x,y)\in F)\land (q\in y)\}\cup\{(x,\emptyset)|(x\in\Dom(F))\land (F(x)=\emptyset) \}
$\\
$
\canonize3(F) = \{(x,s)|\exists y:((x,y)\in F)\land (s\subseteq y) \land (!q:q\in s)\}
$\\

\end{definition}

\begin{definition}(Cardinality)\\
\[
|A|=\Ord\cap\bigcap \{f: | \}
\]
\end{definition}

Here V denote a class of all sets, and Ord denotes the class of all ordinal numbers.
\[
|A|:=\mbox{Ord}\cap\bigcap \{\alpha\in\mbox{Ord}|\exists (f:A\to\alpha):(f\mbox{ injective})\}
\]
Note: here we use the intersection of a class which is defined by 
$
(x\in\bigcap Q)\iff(\forall q\in Q:x\in q)
$, therefore $\bigcap\emptyset = V$.
In this case 
\[
(x\mapsto|x|):V\to\mbox{Ord}
\]
This definition allows also obtain a cardinality of any proper class P, in particular
\[
|P|=\mbox{Ord}
\]
This definition is natural since it agrees with the axiom of limitation of size which implies bijection between V and any proper class.


\begin{definition}(Family)
\[
\]
\end{definition}

\begin{definition}(Family)
\[
(F:A\to\mathcal{U})
\IFF
(F\subseteq V\times V)
\land
(\Dom(F)=A)
\]
\end{definition}
\begin{definition}(Canonical Family)
Every family has a representation which is canonical.
If we use axiomatic approach for families, they should be indistinguishable wrt all important properties.
So agreement may be ocnsidered: "Read canonize(F) in any statement when F is mentioned" or "Add (canonize(F)=F) to the definition".
\end{definition}



\begin{definition}(Family)
A family is a relation which range of relation consists only of subsets of singletones.
\[
(F\family) \IFF (F\subseteq V\times\{-\}[])
\]
\end{definition}
A family is weaker than a funtional symbol, because functional symbols can take proper classes as arguments.\\
A function is weaker than a family because it can return only sets.\\
But the behaviour of $\{-\}$ is well-defined only on sets.
Maybe it should be a family?\\
What about $\{-,-\}$?\\
Should $\{V\}$ be $\emptyset$ or $V$?\\
E.g. Metamath uses an empty set as a default value.\\
So it seems to legit, because it cannot be obtained in any other way out of the unordered pair or ordered pair. At worst it would be a singleton, but never an empty set. So the choice of $\emptyset$ as a default value is justified.\\

It's not the same in other cases, e.g. for "f(x)", therefore it requires an additional assumption to be treated as "legit".

\subsection{Value of a function via set-descriptors}
\begin{definition}
\label{def:value}
Intended use: value of a class-function. Class F is arbitrary.
\[
F(x) = \left(\iota y: \langle x,y\rangle\in F\right)
\]
\end{definition}





\begin{definition}
\label{def:univocalat}
\[
(F\univocalat x)\iff(!y : \langle x,y\rangle \in F)
\]
\end{definition}
\begin{definition}
\label{def:univocal}
\[
(F\univocal)\iff (\forall x: F\univocalat x)
\]
\end{definition}
\begin{theorem}
\label{thm:pair in function}
If $(F\univocalat x)\land(y\set)$, then
\[
\left(F(x)=y\right)\iff\left(\langle x,y\rangle\in F\right)
\]
\begin{align*}
(F\univocalat x)
\implies
\left(
\left(\langle x,y\rangle\in F\right)
\implies
\left(F(x)=y\right)
\right)
\\
(y\set)
\implies
\left(
\left(F(x)=y\right)
\implies
\left(\langle x,y\rangle\in F\right)
\right)
\end{align*}
\iffalse
\[
\mbox{x set!!??}
(F\univocalat x)
\implies
\left(
\left(
\left(F(x)=y\right)\land(y\set)
\right)
\iff\left(\langle x,y\rangle\in F\right)
\right)
\]
\fi
\end{theorem}
\Proof
For formulas to be short, let's define 
\[
Q:=\{y| (\langle x,y\rangle\in F) \land (!y: \langle x,y\rangle\in F)\}.
\]
\noindent($\Leftarrow$) Assume $(\langle x,y\rangle\in F)$.
% USELESS
%\[
%(F\univocal)
%\IMPLIES
%(\forall x:\ !y : \langle x,y\rangle \in F)
%\]
% USELESS
%\[
%(\forall x:\ !y : \langle x,y\rangle \in F)
%\IMPLIES
%(!y : \langle x,y\rangle \in F)
%\]
\[
(F\univocalat x)
\IMPLIES
(!y : \langle x,y\rangle \in F)
\quad\mbox{(by def \ref{def:univocalat})}
\]
\[
(\langle x,y\rangle\in F)
\IMPLIES
(\exists y : \langle x,y\rangle \in F)
\]
\[
(\exists y : \langle x,y\rangle \in F)
\AND
(!y : \langle x,y\rangle \in F)
\IMPLIES
(\exists! y : \langle x,y\rangle \in F)
\]
\[
(\exists! y : \langle x,y\rangle \in F)
\IMPLIES
(\{y | \langle x,y\rangle \in F\}\singleton)
\]
\[
\left(\{y| \langle x,y\rangle\in F\}\singleton\right)
\IMPLIES
\left(\intersection\{y| \langle x,y\rangle\in F\}\right)
\in
\left(\{y| \langle x,y\rangle\in F\}\right)\quad(\mbox{by thm \ref{thm:intersection of singleton} }).
\]
\[
\left(\intersection\{y| \langle x,y\rangle\in F\}\right)
\in
\left(\{y| \langle x,y\rangle\in F\}\right)
\IMPLIES
\langle x,\left(\intersection\{y| \langle x,y\rangle\in F\}\right)\rangle\in F
\]
\[
(!y : \langle x,y\rangle \in F)
\AND
(\langle x,y\rangle\in F)
\AND
\langle x,\left(\intersection\{y| \langle x,y\rangle\in F\}\right)\rangle\in F
\IMPLIES
\left(\intersection\{y| \langle x,y\rangle\in F\}\right)=y
\]

%\[
%\beta
%\IMPLIES
%\left(
%\alpha
%\iff
%(\alpha \land \beta))
%\right)
%\]
\[
\beta\implies(\alpha\iff(\alpha \land \beta)))
\]
\[
(!y: \langle x,y\rangle\in F)
\IMPLIES
\left(
(\langle x,y\rangle\in F)
\iff
(\langle x,y\rangle\in F \land (!y: \langle x,y\rangle\in F))
\right)
\]

%(!y: \langle x,y\rangle\in F)
\[
\IMPLIES
\left(
\{y| \langle x,y\rangle\in F\}
=
Q
\right)
\]


%\[
%\left(\{y| \langle x,y\rangle\in F\}
%=
%\{y| \langle x,y\rangle\in F \land (!y: \langle x,y\rangle\in F)\}\right)
%\IMPLIES
%\]
\[
\IMPLIES
\left(\intersection\{y| \langle x,y\rangle\in F\}\right)=
\left(\intersection
Q
\right)
\]

\iffalse
%\{y| \langle x,y\rangle\in F \land (!y: \langle x,y\rangle\in F)\}
\[
\left(\{y| \langle x,y\rangle\in F\}
=
Q
\right)
\AND
\left(\intersection\{y| \langle x,y\rangle\in F\}\right)=y
\IMPLIES
\left(
\intersection Q
\right)
=y
\]

\[
\left(\intersection
Q
\right)=y
\AND
(\iota y: (\langle x,y\rangle\in F))
=
\left(\intersection
Q
\right)
\IMPLIES
\left((\iota y: \langle x,y\rangle\in F)=y\right)
\]
\[
\left((\iota y: \langle x,y\rangle\in F)=y\right)
\IMPLIES
(F(x)=y)
\]
So we proved that F(x)=y.
\fi

\begin{align*}
F(x)&=\left((\iota y: \langle x,y\rangle\in F\right)\\
&=\intersection Q\\
&=\intersection \{y: \langle x,y\rangle\in F\land (!y:\langle x,y\rangle\in F)\}\\
&=y
\end{align*}
So we proved that $(F(x)=y)$.\\
\noindent($\Rightarrow$) Assume $(F(x)=y)$. To show:$(\langle x,y\rangle\in F)$
\begin{align*}
y&=F(x)
\quad(\mbox{by assumption}) \\
&=\left(\iota y: \langle x,y\rangle\in F\right)
\quad(\mbox{by def of value of a function \ref{def:value}})
\\
&=\left(\intersection Q\right)
\quad(\mbox{by def of Q and def of a iota operator \ref{def:iota operator}})
\end{align*}
%\[
%(F(x)=y)
%\IMPLIES
%\left(\iota y: \langle x,y\rangle\in F\right)=y.
%\]
%\[
%\left(\iota y: \langle x,y\rangle\in F\right)=y
%\IMPLIES
%\left(\intersection Q\right)=y.
%\]
\[
(F\univocal)
\IMPLIES
(\forall x:\ !y : \langle x,y\rangle \in F)
\]

\[
(\forall x:\ !y : \langle x,y\rangle \in F)
\IMPLIES
(!y : \langle x,y\rangle \in F)
\]

\[
(y\set)
\AND
\left(\intersection Q\right)=y
\IMPLIES
\left(\left(\intersection Q\right)\set\right)
\]

\[
\left(\left(\intersection Q\right)\set\right)
\IMPLIES
(\exists w: w\in Q)
\quad(\mbox{by thm \ref{thm:intersection set then}}).
\]

%GOOD, BUT USELESS
%\[
%(w\in Q)
%\IMPLIES
%\langle x,w\rangle\in F
%\land 
%(!w: \langle x,w\rangle\in F)
%\quad(\mbox{by def of Q}).
%\]

%(!w:w\in \{y|\varphi\land (!y:\varphi)\})
%(!w:w\in \{y|(\langle x,y\rangle\in F)\land (!y:(\langle x,y\rangle\in F))\})
\[
(!w:w\in Q)
\quad\mbox{(by thm \ref{thm:class with at most one object} with $\varphi:=(\langle x,y\rangle\in F)$)}
\]
%...

\[
(\exists w: w\in Q)
\AND
(! w: w\in Q)
\IMPLIES
(\exists! w: w\in Q)
\]

\[
(\exists! w: w\in Q)
\IMPLIES
(Q\singleton)
\]

\[
\left(Q\singleton\right)
\IMPLIES
\left(\intersection Q\right)\in Q
\quad(\mbox{by thm \ref{thm:intersection of singleton} }).
\]

\[
\left(\intersection Q\right)\in Q
\AND
\left(\intersection Q\right)=y
\IMPLIES
y\in Q.
\]

\[
y\in Q
\IMPLIES
(\langle x,y\rangle\in F)
\land 
(!y: \langle x,y\rangle\in F)
\quad(\mbox{by def of Q}).
\]
So $(\langle x,y\rangle\in F)$ is proven as intended.\\
\Done
\begin{theorem}
\[
((\iota y:\phi)\ldefined)
\implies
\phi[(\iota y:\phi)/y]
\]
\end{theorem}
\Proof
\[
((\iota y:\phi)\ldefined)
\IFF
((\iota y:\phi)\set)
\comment{by def that is valid only for terms denoting sets}
\]

\[
((\iota y:\phi)\set)
\IMPLIES
\left(\intersection\{y:\phi\land(!y:\phi)\}\set\right)
\comment{by def of iota}
\]
\Done

\begin{theorem}
\label{thm:pair in function}
If $(F\univocalat x)\land(y\set)$, then
\[
\left(F(x)=y\right)\iff\left(\langle x,y\rangle\in F\right)
\]
\begin{align*}
1) & & (y\set)
\implies
\left(
\left(F(x)=y\right)
\implies
\left(\langle x,y\rangle\in F\right)
\right)
\\
2) & & (F\univocalat x)
\implies
\left(
\left(\langle x,y\rangle\in F\right)
\implies
\left(F(x)=y\right)
\right)
\\
\end{align*}
\end{theorem}
Alternative proof:\\
\Proof
\noindent($\Rightarrow$)
\begin{align*}
y&=F(x)
\quad(\mbox{by assumption}) \\
&=\left(\iota y: \langle x,y\rangle\in F\right)
\quad(\mbox{by def of value of a function \ref{def:value}})
\\
\end{align*}
%&=\left(\intersection Q\right)
%\quad(\mbox{by def of Q and def of a iota operator \ref{def:iota operator}})
\[
(y=\left(\iota y: \langle x,y\rangle\in F\right))
\AND
(y\set)
\IMPLIES
(\left(\iota y: \langle x,y\rangle\in F\right) \set)
\]
\[
(\left(\iota y: \langle x,y\rangle\in F\right) \set)
\IMPLIES
\left(\exists!(y\set):\langle x,y\rangle\in F\right)
\comment{by thm \ref{thm:iota set exists unique set}}
\]
...\\
\noindent($\Leftarrow$)\\
...\\
\Done
\subsection{Descriptors OLD}
\begin{axiom}
\[
\exists (x\set):\varphi
\implies
(\varepsilon x: \varphi) \in \{x|\varphi\}
\]
\end{axiom}
\begin{axiom}

\[
(\varepsilon x: \varphi) \set
\implies
\exists (x\set):\varphi
\]
\end{axiom}

\begin{axiom}
\[
\exists! (x\set):\varphi
\implies
(\iota x: \varphi) \in \{x|\varphi\}
\]
\end{axiom}
\begin{axiom}
%(\iota x: \varphi) \in \{x|\varphi\}
\[
(\iota x: \varphi) \set
\implies
\exists! (x\set):\varphi
\]
\end{axiom}



It would be nice to define iota operator both for sets and proper classes, but it is not possible in a classical setting.
\[
(I y: \varphi) = \intersection \{y|\varphi \land (!y:\varphi)\}
\comment{bad ``definition''}
\]
In general, the problem of NBG and MK is the principal inability to collect classes. Maybe there is a good way to define iota operator "up to universe", but here it will not be researched.
(the definition is bad, since $\varphi$ may contain a disjunction with the property of any proper class(e.g. V), so the uniqueness will not hold, since $\varphi$ will be true both for a set and for a proper class.)
%So the following definition does not work correctly for proper classes:
\begin{definition}
\label{def:iota operator}
Iota operator for sets.
\[
(\iota y: \varphi) = \intersection \{y|\varphi \land (!(y\set):\varphi)\}
\]
\end{definition}



\begin{theorem}
\label{thm:iota set exists unique set}
%Note: it should be possible to substitute some fresh variable $w$ instead of $\varphi$.
%This may not be true:
%\[ ???
%((\iota y: \varphi) \set)
%\iff
%(\exists!(b\set):\varphi)
%\]
%Variant2(true before redef but ugly):
%\[
%((\iota y: \varphi) \set)
%\iff
%(\exists(b\set):\varphi)
%\land
%(!b:\varphi)
%\]
%Variant3:
%\[
%((\iota y: \varphi) \set)
%\iff
%(\exists!b:\varphi)
%\]
%Variant4:
%\[
%((\iota (y\set): \varphi) \set)
%\iff
%(\exists!(y\set):\varphi)
%\]
%Variant5:
\[
((\iota y: \varphi) \set)
\iff
(\exists!(y\set):\varphi)
\]

\end{theorem}
\Proof
%For formulas to be short, let's locally define the term Q:
%\[
%Q:=\{y|\varphi \land (!(y set):\varphi)\}
%\]
\noindent$(\Rightarrow)$\\
Assume $((\iota y: \varphi) \set)$. To show $(\exists!(y\set):\varphi)$.
\[
((\iota y: \varphi) \set)
\IMPLIES
\left(\intersection \{y|\varphi \land (!(y\set):\varphi)\} \set\right)
\comment{by def \ref{def:iota operator}}
\]
\[
\left(\intersection \{y|\varphi \land (!(y\set):\varphi)\} \set\right)
\IMPLIES
\left(\{y|\varphi \land (!(y\set):\varphi)\}\nonempty\right)
\]
\[
\left(\{y|\varphi \land (!(y\set):\varphi)\}\nonempty\right)
\IMPLIES
\left(\exists	y: y\in \{y|\varphi \land (!(y\set):\varphi)\}\right)
\]
\[
\left(
y\in \{y|\varphi \land (!(y\set):\varphi)\}
\right)
\IMPLIES
\left(
(y\set)\land \varphi\land (!(y\set):\varphi)
\right)
\]
\[
(y\set)
\AND
\varphi
\IMPLIES
\left(
\exists (y\set):\varphi
\right)
\]
%\[
%\left(
%!y:\varphi
%\right)
%\IMPLIES
%\left(
%!(y\set):\varphi
%\right)
%\]
\[
\left(
\exists (y\set):\varphi
\right)
\AND
\left(
!(y\set):\varphi
\right)
\IFF
\left(
\exists! (y\set):\varphi
\right)
\]

\noindent$(\Leftarrow)$\\
Assume $(\exists!(y\set):\phi)$. To show:$((\iota y:\phi)\set)$.
%%Here $!x:$.
\[
(\exists!(y\set):\phi)
\IMPLIES
(\exists(y\set):\phi)
\land
(!(y\set):\phi)
\]
\[
(y\set)
\AND
\phi
\AND
(!(y\set):\phi)
\IMPLIES
y\in\{y:\phi\land(!(y\set):\phi)\}
\]
\[
y\in\{y:\phi\land(!(y\set):\phi)\}
\IMPLIES
(\{y:\phi\land(!(y\set):\phi)\}\nonempty)
\]
\[
(\{y|\varphi \land (!y:\varphi)\}\nonempty)
\IMPLIES
\left(\intersection \{y|\varphi \land (!y:\varphi)\} \set\right)
\]
\[
\left(\intersection \{y|\varphi \land (!y:\varphi)\} \set\right)
\IMPLIES
\left(\left(\iota y: \varphi\right) \set\right)
\]
\Done
